% Выбираем класс extarticle для поддержки шрифта 14pt
\documentclass[14pt, a4paper]{extarticle}

% ================= МОВА ТА ШРИФТИ =================
\usepackage[utf8]{inputenc}
\usepackage[T2A]{fontenc}
\usepackage[ukrainian]{babel}

% Times New Roman для тексту (з правильною кирилицею)
\usepackage{tempora} 

% ================= ГЕОМЕТРІЯ ТА ПОЛЯ =================
% Зліва 30мм (для зшивання), інші по стандартам з тексту
\usepackage[left=3cm, right=1cm, top=2cm, bottom=2cm]{geometry}

% ================= АБЗАЦИ ТА ІНТЕРВАЛИ =================
\usepackage{setspace}
\onehalfspacing % Міжрядковий інтервал 1.5

\usepackage{indentfirst} % Відступ для першого абзацу в розділі
\setlength{\parindent}{1.25cm} % Абзацний відступ рівно 1.25 см

% ================= ОФОРМЛЕННЯ ЗАГОЛОВКІВ =================
\usepackage{titlesec}

% Розділ 1 (РОЗДІЛ 1, з нового рядка НАЗВА ВЕЛИКИМИ ЛІТЕРАМИ, по центру)
\titleformat{\section}[display]
{\filcenter\bfseries} % По центру, напівжирний
{\MakeUppercase{Розділ \thesection}} % Слово РОЗДІЛ X
{0pt} % Відстань між словом РОЗДІЛ і назвою
{\MakeUppercase} % Назва великими літерами

% Щоб кожен розділ починався з нової сторінки
\newcommand{\sectionbreak}{\clearpage}

% Відстань від назви розділу до тексту (2 рядки)
\titlespacing*{\section}{0pt}{*0}{2\baselineskip}

% Підрозділ 1.1 (З абзацу, напівжирний, без крапки в кінці)
\titleformat{\subsection}[block]
{\bfseries}
{\hspace{1.25cm}\thesubsection} % Відступ номера як абзац
{1ex}
{}

% Відстань для підрозділів (1 рядок до і після)
\titlespacing*{\subsection}{0pt}{\baselineskip}{\baselineskip}

% Пункт 1.1.1 (Звичайний шрифт, з абзацу)
\titleformat{\subsubsection}[block]
{\normalfont}
{\hspace{1.25cm}\thesubsubsection}
{1ex}
{}

\titlespacing*{\subsubsection}{0pt}{\baselineskip}{\baselineskip}

% ================= ФОРМУЛИ ТА РИСУНКИ =================
\usepackage{amsmath}
\usepackage{amsthm}
% Times New Roman для математичних формул
\usepackage{newtxmath} 
\newtheorem{statement}{Твердження}
\newtheorem{theorem}{Теорема}
\theoremstyle{definition} 
\newtheorem{remark}{Зауваження}
% Відступи до і після формул (12 пт)
\setlength{\abovedisplayskip}{12pt}
\setlength{\belowdisplayskip}{12pt}
\setlength{\abovedisplayshortskip}{12pt}
\setlength{\belowdisplayshortskip}{12pt}

% Для рисунків (відступ 12 пт)
\usepackage{caption}
\captionsetup{aboveskip=12pt, belowskip=12pt}
\setlength{\intextsep}{12pt} % Відстань від тексту до рисунка

% ================= НУМЕРАЦІЯ СТОРІНОК =================
\usepackage{fancyhdr}
\pagestyle{fancy}
\fancyhf{} % Очищаємо всі колонтитули
\fancyfoot[R]{\thepage} % Номер сторінки знизу справа
\renewcommand{\headrulewidth}{0pt} % Прибираємо лінію зверху
\renewcommand{\footrulewidth}{0pt} % Прибираємо лінію знизу

% ================= ІНШІ ПАКЕТИ =================
\usepackage{graphicx}
\usepackage{float}
\usepackage{xurl}
\usepackage{csquotes} % Для правильних лапок (працює з babel)

\sloppy % Дозволяє розтягувати пробіли, щоб текст не вилазив за поля
